\chapter{Phantom Pain (Phantom Limb Pain)}

\section*{Definition}
Phantom pain is the sensation of pain or discomfort coming from a body part that is no longer there, most commonly after amputation of a limb. It is different from stump pain (pain at the surgical site). The sensation can feel like the missing limb is still present and experiencing pain, burning, cramping, or other unpleasant feelings.

\section*{What Happens in Your Body}
After amputation, the nerves that previously supplied the removed limb continue to send signals to the spinal cord and brain. The brain's sensory map still has a representation of the missing limb, and it interprets signals from the severed nerves as coming from the limb that is no longer there. This is called neuroplastic remapping gone awry — the brain has not yet adjusted to the loss of the limb. The spinal cord can also become hypersensitive, amplifying pain signals. The result is a vivid, often painful sensation that feels as though it originates in the missing body part.

\section*{Causes}
\begin{itemize}
  \item Limb amputation (most common) — can occur immediately after surgery or weeks, months, or years later
  \item Spinal cord injury
  \item Nerve avulsion (nerve torn from the spinal cord)
  \item Removal of a breast (phantom breast pain after mastectomy)
  \item Removal of an eye (phantom eye pain after enucleation)
  \item Removal of internal organs (phantom rectum or bladder pain after surgical removal)
\end{itemize}

\section*{When to See a Doctor}
See your doctor if:
\begin{itemize}
  \item Phantom pain is severe, persistent, or interfering with your daily life
  \item The pain is not controlled by current treatments
  \item You have signs of depression, anxiety, or sleep disturbance related to the pain
  \item You notice new redness, swelling, or discharge at the stump site (possible infection)
  \item The nature of the pain changes suddenly or dramatically
\end{itemize}

\section*{Related Symptoms}
\begin{itemize}
  \item Stump pain (pain at the surgical site, different from phantom pain)
  \item Sensation that the missing limb is still present (phantom limb sensation, not necessarily painful)
  \item Tingling, electric shock, or burning sensations in the missing limb
  \item Cramping or squeezing sensations
  \item Feeling that the missing limb is in an abnormal position (telescoping — sensation that the limb is shortening)
  \item Depression, anxiety, or grief related to the amputation
  \item Sleep disturbance
\end{itemize}
\textbf{Important:} Phantom pain is a real medical condition, not a psychological problem. It results from changes in the nervous system and requires proper treatment.

\section*{Self-Care / Home Management}
\begin{itemize}
  \item Mirror therapy — using a mirror to create a reflection of the intact limb so it appears the missing limb is present and moving; this can help reduce pain over time.
  \item Distraction techniques — engaging in activities that draw your attention away from the pain.
  \item Relaxation techniques such as deep breathing, meditation, or guided imagery.
  \item Desensitization — gently touching or tapping the stump with different textures (soft cloth, cotton ball) can help retrain nerve signals.
  \item Exercise and physical therapy to maintain strength and mobility.
  \item Wearing a well-fitted prosthetic limb can reduce phantom pain in some people.
  \item Connect with support groups for people with limb loss.
\end{itemize}

\section*{How Your Doctor Will Evaluate You}
\begin{itemize}
  \item Your doctor will ask about the location, quality, and timing of the pain, and whether it differs from any pain at the surgical site.
  \item They will examine the stump for signs of infection, poor healing, or neuroma (a thickened nerve end that can cause pain).
  \item They will assess your mental health, sleep quality, and overall function.
  \item Imaging (X-ray, ultrasound, or MRI) of the stump may be done to check for bone spurs, neuromas, or other structural problems.
  \item Referral to a pain specialist or a physiatrist (rehabilitation doctor) is common.
\end{itemize}

\section*{Treatment Options}
\begin{itemize}
  \item Medications: Anticonvulsants (gabapentin, pregabalin), antidepressants (amitriptyline, nortriptyline), opioids (for severe pain, short-term), and NSAIDs for associated inflammation.
  \item Mirror therapy — strong evidence supports its effectiveness; a therapist can teach you the technique.
  \item Nerve blocks — injecting anesthetic near the nerve can provide temporary relief.
  \item Spinal cord stimulation — an implanted device sends electrical pulses to the spinal cord to interrupt pain signals.
  \item Transcutaneous electrical nerve stimulation (TENS) — a device that delivers mild electrical current through the skin to reduce pain.
  \item Cognitive behavioral therapy (CBT) to help manage the psychological impact of chronic pain and amputation.
  \item Acupuncture and biofeedback have helped some people.
  \item Virtual reality therapy — newer approach that uses immersive environments to trick the brain into seeing the missing limb move.
  \item Repetitive transcranial magnetic stimulation (rTMS) — experimental but promising.
\end{itemize}

\section*{References}
\begin{thebibliography}{99}
\bibitem{ref1} Nikolajsen L, Jensen TS. "Phantom limb pain." \textit{British Journal of Anaesthesia}, 2001;87(1):107--116.
\bibitem{ref2} Flor H. "Phantom-limb pain: characteristics, causes, and treatment." \textit{Lancet Neurology}, 2002;1(3):182--189.
\bibitem{ref3} Ramachandran VS, Rogers-Ramachandran D. "Synaesthesia in phantom limbs induced with mirrors." \textit{Proceedings of the Royal Society of London B}, 1996;263(1369):377--386.
\bibitem{ref4} Hsu E, Cohen SP. "Postamputation pain: epidemiology, mechanisms, and treatment." \textit{Journal of Pain Research}, 2013;6:121--136.
\bibitem{ref5} Weeks SR, Anderson-Barnes VC, Tsao JW. "Phantom limb pain: theories and therapies." \textit{Neurologist}, 2010;16(5):277--286.
\bibitem{ref6} Flor H, Diers M, Andoh J. "The neural basis of phantom limb pain." \textit{Trends in Cognitive Sciences}, 2013;17(7):307--308.

\end{thebibliography}

\clearpage
